


\section{\ours{}: Turning Binary Rewards into Dense Self-Supervision} \label{sec:method}

% \danqi{Title: reward or rewards}

%\yj{Also, how about having a clear summary of the full method with bullet points? - Phase1:... - Phase 2: ...}
We consider a dataset $\mathcal{D} = \{(x_i, a_i)\}_{i=1}^N$, where $x_i$ denotes an input problem and $a_i$ its corresponding ground-truth final answer. Crucially, we do not assume access to gold solutions. Let $\pi_\theta$ denote the student policy with parameters $\theta$, which generates a reasoning response $y$ for input $x$ according to $\pi_\theta(y \mid x)$. For each example $(x, a)$ and generated response $y$, we assume access to a binary reward $r(y,a) \in \{0,1\}$, where $r(y,a)=1$ if the final answer extracted from $y$ matches $a$, and $r(y,a)=0$ otherwise.

The core design of \ours{} is around a single model that can play two roles: a \emph{generator} and a \emph{reviser}. We partition the training set into two disjoint subsets of sizes $N_1$ and $N_2$:
\begin{itemize}
\setlength{\itemsep}{-0.2em}
\item  \textbf{Phase 1, Self-Revision Training}: Using the first subset of $N_1$ examples, we train the model to perform both the generator and reviser roles.\newline
\item   \textbf{Phase 2, \opsd{}}: Using the remaining $N_2$ examples, we leverage the reviser to transform sparse outcome-level supervision into dense token-level supervision over the generator's responses.
\end{itemize}

In our experiments, we refer to the model obtained after Phase 1 as the \sft{} model, and the model obtained after Phase 2 as the \ours{} model.

%In the first stage, \sft{}, we use $N_1$ examples to train the model to perform the reviser role. We then define our algorithm \ours{} on the remaining $N_2$ examples, where the reviser transforms sparse outcome-level supervision into dense token-level feedback for the generator. \yh{mention ablations on N1 and N2}

% \subsection{Preliminaries}
% \yj{Maybe it's better to have a preliminary subsection that outlines the notations, SFT, and OPD?}\\

%\ap{iffy - need to modify} \ours{} treats self-improvement as a dual-role process in which a single model alternates between a \emph{generator} and a \emph{reviser}. The generator generates candidate reasoning responses from $\pi_\theta(\cdot \mid x)$. The reviser then conditions on the sampled trajectory and its outcome reward: if the trajectory is incorrect, it is prompted to revise it toward a better solution; if it is correct, it is prompted to rephrase it into an alternative valid trajectory. The resulting revised responses provide token-level supervision on the generator’s own rollouts. Empirically, however, we find that off-the-shelf reasoning models are much stronger generators than revisers \ap{cite the relevant section/table}. We therefore first perform a revision warm-start stage that improves the model’s self-revision ability while preserving its generation quality. We then run on-policy distillation, using the reviser’s outputs as token-level targets for training the generator.

% We consider a dataset $\mathcal{D} = \{(x_i, a_i)\}_{i=1}^N$, where $x_i$ is an input problem and $a_i$ is the corresponding ground-truth final answer. Importantly, we do not assume access to ground-truth reasoning traces. We will use $\pi_{\theta}$ to denote the student model with parameter $\theta$, which generates a reasoning response $y$ given input $x$ from its policy $\pi_\theta(y \mid x)$. Given a response $y$ for an input $(x, a)$, we assume access to a reward model $r(y, a) \in \{0, 1\}$ that returns 1 if the final answer extracted from $y$ is equal to $a$.

% We require a model to show 2 roles for \ours{}: a generator and a reviser. We find that off-the-shelf models are good generators for reasoning tasks, but struggle to revise their responses when wrong \ap{Need to refer to the right table/section here}. So, we have an initial phase of training, where we train to reliably perform self-revision, while also retaining its performance to generate good responses. Then, we train the generator using an on-policy distillation strategy, where the reviser returns per token supervision on the responses of the generator.


% In \ours{}, we train the model first to play two roles; a generator and a reviser. 


% \subsubsection{Step 1: Critique Training}

% Primary steps for generating data:

% Idea:
%  - Generate multiple responses on prompts from the given dataset.
%  - On responses where the model is wrong, we ask the model to generate a critique and refine the initial response.
%  - On responses where the model is right, we simply ask the model to rephrase the response.
%  - Create critique dataset of the form - (x, y+, rephrased y) and (x, y-, critiqued y)

% \paragraph{Setup.}
% We first train the model to revise its own sampled responses using only access to the ground-truth final answer, without requiring ground-truth reasoning traces. Let $\mathcal{D} = \{(x_i, a_i)\}_{i=1}^N$ denote the training set, where $x_i$ is an input problem and $a_i$ is the corresponding ground-truth final answer. We emphasize that we do not assume access to any gold chain-of-thought or reasoning trace.


\subsection{Phase 1 (\sft{}): Self-Revision Training} \label{sec:phase1}

\looseness-1Our base models exhibit strong \textit{generator} performance but weak \textit{reviser} performance (\Cref{fig:self-revision}). This phase strengthens the model as a \textit{reviser}, while also improving it as a \textit{generator}.
%For each problem, we sample multiple responses and evaluate them with binary reward $r$. For incorrect responses, we prompt the model to correct its response; for correct responses, we prompt it to rephrase. This produces a self-generated revision dataset that trains the model to condition its behavior on the reward of its initial response. More formally, the procedure is as follows.
% \looseness-1Our base models show strong \textit{generator} performance but weak \textit{reviser} performance (see \Cref{fig:self-revision}). This phase aims to strengthen the  performance of the model as a \textit{reviser}, while  improving its performance as a \textit{generator} as well.

% \looseness-1We sample multiple responses from the model for each problem and evaluate them using the binary reward $r$. For responses that are incorrect, we prompt the model to revise and produce a new response; for responses that are correct, we prompt it to generate a rephrased version of the response. This yields a self-generated revision dataset that teaches the model to respond differently depending on the reward of its initial response. More formally, we have the following steps.

%After training on this dataset, we show that (a) model's revised responses on incorrect responses improves substantially, and (b) model's rephrased response is shorter compared to the length of incorrect response.

%\jiarui{Did we explain our motivation for getting rephrased correct responses?} \ap{Need to: asking yinghui} 
 
% \paragraph{On-policy response generation.}

% For each input $x$, we sample multiple initial reasoning responses
% \[
% y \sim \pi_\theta(\cdot \mid x).
% \]
%We use a binary verifier $r(y, a) \in \{0,1\}$ to determine whether the final answer extracted from $y$ matches the ground-truth answer $a$.
%\yh{I replaced all $a$ with $r$ because we're not telling the model the ground truth answer. Yet we need to better define $r$ and $r(y, a)$}
%\yh{I feel the $p$ notation is a bit confusing here}

\paragraph{Outcome-Conditioned Self-Revision.}
For each input $x$, we sample multiple initial reasoning responses $y_\text{init} \sim \pi_\theta(\cdot \mid x)$. Given a sampled response $y_\text{init}$, we construct a revision prompt that conditions on the original problem, the sampled response, and a control phrase indicating whether the response should be revised or simply rephrased: 
\[
P_{r} =
\begin{cases}
\text{``Let me rephrase the above solution.''}, & \text{if } r(y_\text{init},a)=1, \\
\text{``Wait, this response is not correct, let me start over.''}, & \text{if }  r(y_\text{init},a)=0.
\end{cases}
\]
We then use the same model to generate a revised response:
\[
y_\text{revised} \sim \pi_\theta(\cdot \mid x, y_\text{init}, P_{r}).
\]

For correct initial responses, this encourages the model to produce a rephrase of a correct solution, for incorrect initial responses, this encourages the model to critique on the previous attempt and regenerate a correct solution. We observe that when asked to rephrase a correct response, the model typically produces a shorter answer, which may additionally encourage more concise responses during the second phase. We defer further discussion to \Cref{sec:experiments}. We create a dataset $\mathcal{D}_{\textsc{revision}} =  \{(x,y_\text{init},P_r, y_\text{revised})\}$.

% \[
% \Tilde{y} \sim \pi_\theta(\cdot \mid x, y, P_{r}).
% \]
%\yj{Again, $z$ or $y_{revised}$?} 
%\sk{I like $y_{revised}$ a bit better, I think it better communicates the meaning}
% \yj{Also, "Rephrase" needs more justification as it introduces extra compute cost and the original y is also sampled from the student model}\jiarui{+1 I feel like we need to explain why we do rephrasing}
 
% \paragraph{Self-revision dataset construction.}
% We create two datasets as follows:
% \begin{align*}
% \mathcal{D}_{\textsc{revision}} &= \{(x,y,P_r, \Tilde{y}) \,:\, r(\Tilde{y}, a)=1, r(A, a) = 0\}, \\
% \mathcal{D}_{\textsc{rephrase}} &= \{(x,y,P_r, \Tilde{y}) \,:\, r(\Tilde{y}, a)=1, r(A, a) = 1\}
% \end{align*}
% \yh{hide notation for D-rephrase}


%, which has the following two components:

%\textbf{Self-revision objective.}
% We then train the student on the generated datasets with loss $\mathcal{L}_{\sft{}}(\theta)$: 
%Now, we train the model with the following loss $\mathcal{L}_{\text{\sft{}}}$ consisting of two log-likelihood terms $\mathcal{L}_{\text{generation}}$ and $\mathcal{L}_{\text{revision}}$:
% on $(y_\text{init},P_r, y_\text{revised})$, conditioned on the input $x$ in the dataset $D_{\textsc{revision}}$: 
%on a combination of two objectives:

% yj: I think this looks a bit visually dense right now. It may read better if the revision loss, generation loss, and final objective are written as separate numbered equations rather than in a single small-font aligned block. That would make each term easier to parse and easier to reference later in the paper. => Made an edit
\textbf{Self-Revision Objective.} We train the model on two tasks simultaneously: given the input, initial attempt, and reward signal, produce a corrected or rephrased response (revision); and given only the input, produce the full correct response from scratch (generation). The corresponding loss $\mathcal{L}_{\text{\sft{}}}$ consists of two log-likelihood terms $\mathcal{L}_{\text{revision}}$ and $\mathcal{L}_{\text{generation}}$.

The \emph{revision loss} trains the model to produce $y_\text{revised}$ when conditioned on the input $x$, the first attempt $y_\text{init}$, and the reward prompt $P_r$:
{\small
\begin{equation}
\mathcal{L}_{\text{revision}}(\theta)=
\mathbb{E}_{(x,y_\text{init},P_r, y_\text{revised})\sim \mathcal{D}_{\textsc{revision}}}
\left[
-\sum_{t=1}^{|y'|}
\log \pi_\theta\!\left(y'_t \mid x, y_\text{init}, P_r, y'_{<t}\right)
\right],\text{where } y' = y_\text{revised}.
\label{eq:revision}
\end{equation}
}
The \emph{generation loss} retains the model's ability to produce the full correct response conditioned only on input $x$:
{
% \small
\fontsize{8.8pt}{9pt}\selectfont
\begin{equation}
\mathcal{L}_{\text{generation}}(\theta)=
\mathbb{E}_{(x,y_\text{init},P_r, y_\text{revised})\sim \mathcal{D}_{\textsc{revision}}}
\left[
-\sum_{t=1}^{|y'|}
\log \pi_\theta\!\left(y'_t \mid x, y'_{<t}\right)
\right],\text{where } y' = [y_\text{init},P_r, y_\text{revised}].
\label{eq:generation}
\end{equation}
}
The overall self-revision objective combines these two terms:
{
\begin{equation}
\mathcal{L}_{\text{\sft{}}}(\theta)=
\mathcal{L}_{\text{revision}}(\theta) + \mathcal{L}_{\text{generation}}(\theta).
\label{eq:sft}
\end{equation}}

\looseness-1 $\mathcal{L}_{\text{revision}}$ (Eq.~\ref{eq:revision}) trains the model to self-revise, i.e., producing $y_\text{revised}$ conditioned on the input $x$, the first attempt $y_\text{init}$, and the prompt $P_r$, while $\mathcal{L}_{\text{generation}}$ (Eq.~\ref{eq:generation}) retains the model's generation capability, i.e., producing correct response $y$ conditioned on input $x$. This combination implicitly teaches the model to actively evaluate its current response at inference time and generate a self-revision if necessary. As a result, the trained model often exhibits explicit self-revision behaviors that lead to extremely long responses (see \Cref{sec:srt-results,sec:distillation-results}). In the next phase, we perform on-policy distillation on the \sft{} model, distilling this learned revision behavior back into the generator to enable sample-efficient self-improvement.

%%%% Backup
% \textbf{Self-revision objective.} We train the model on two tasks simultaneously: given the input, initial attempt, and reward signal, produce a corrected or rephrased response (revision); and given only the input, produce the full correct response from scratch (generation). The corresponding loss $\mathcal{L}_{\text{\sft{}}}$ consisting of two log-likelihood terms $\mathcal{L}_{\text{revision}}$ and $\mathcal{L}_{\text{generation}}$:

% {\small
% \begin{align*}
% \mathcal{L}_{\text{revision}}(\theta)=&
% \mathbb{E}_{(x,y_\text{init},P_r, y_\text{revised})\sim \mathcal{D}_{\textsc{revision}}}
% \left[
% -\sum_{t=1}^{|y'|}
% \log \pi_\theta\!\left(y'_t \mid x, y_\text{init}, P_r, y'_{<t}\right)
% \right] \quad \quad \text{where } y' = y_\text{revised}.\\
% \mathcal{L}_{\text{generation}}(\theta)=&
% \mathbb{E}_{(x,y_\text{init},P_r, y_\text{revised})\sim \mathcal{D}_{\textsc{revision}}}
% \left[
% -\sum_{t=1}^{|y'|}
% \log \pi_\theta\!\left(y'_t \mid x, y'_{<t}\right)
% \right] \quad \quad \text{where } y' = [y_\text{init},P_r, y_\text{revised}].\\
% \mathcal{L}_{\text{\sft{}}}(\theta)=&
% \mathcal{L}_{\text{revision}}(\theta) + \mathcal{L}_{\text{generation}}(\theta).\\
% \end{align*}
% }
% \[
% \mathcal{L}_{\text{\sft{}}}(\theta)=
% \mathcal{L}_{\text{revision}}(\theta) + \mathcal{L}_{\text{generation}}(\theta).\\
% \]
% &+
% \mathbb{E}_{(Q,A,P_r,\tilde{A})\sim \mathcal{D}_{\textsc{revision}} \cup \mathcal{D}_{\textsc{rephrase}}}
% \left[
% -\sum_{t=1}^{|\tilde{A}|}
% \log \pi_\theta\!\left(\tilde{A}_t \mid Q, A, P_r, \tilde{A}_{<t}\right)
% \right],


% where $y = [A,\, P_r,\, \tilde{A}]$, the concatenation of the initial attempt $A$, the control phrase $P_r$, and the revised attempt $\Tilde{A}$.

% \looseness-1 The first objective trains the model to predict $P_r$ from its initial attempt. This encourages the model to implicitly reason about whether its first response is flawed and to critique it before producing a revised second attempt. We apply this objective asymmetrically only on $D_{\textsc{\textsc{rephrase}}}$, where the model must correct an incorrect response, as we did not observe benefits from applying it on $D_{\textsc{rephrase}}$. This asymmetric training is also the primary reason that, by the end of this phase, the model tends to produce two attempts during inference, with the second attempt revising the first.
 
% \looseness-1 The second term trains the model to \emph{adapt its behavior} based on the reward verdict of its initial attempt. Conditioned on $P_r$, the model learns to produce a \emph{correction} when the original response is incorrect, and a \emph{rephrase} when it is correct. 


%We primarily train the model to criticize its first attempt in order to generate the second attempt.

% \looseness-1 The first term trains the model to \emph{self-correct}. On examples in $\mathcal{D}_{\textsc{revision}}$, the model is optimized to generate both $P_{y,r}$ and $z$; that is, it is trained to first recognize that its initial response is incorrect, and then produce a corrected second attempt $z$.

% \looseness-1The second term trains the model to \emph{differentiate its behavior based on the prompt}. Conditioned on the prompt $P_{y,r}$, which contains the original response $y$ together with its binary outcome $r$, the model is trained to generate $z$: a \emph{revision} when the original response is incorrect, and a \emph{rephrase} when it is correct. This term explicitly teaches the model to condition its behavior on the reward in the prompt, i.e., to revise when incorrect and rephrase when correct.

%In summary, this phase trains the model to act as a reviser whose behavior is conditioned on the prompt $P_r$, which contains the scalar reward. However, as one can observe, the model is also trained to implicitly determine whether its initial response is correct or incorrect, and then generate a revised second attempt $\Tilde{y}$. 

% \looseness-1 The self-revision objective $\mathcal{L}_{\text{\sft{}}}$ serves two purposes: (1) $\mathcal{L}_{\text{revision}}$ trains the model to self-revise, i.e., producing $y_\text{revised}$ when conditioned on the input $x$, the first attempt $y_\text{init}$, and the prompt $P_r$, (2) $\mathcal{L}_{\text{generation}}$ retains the model's generation capability, i.e., producing correct response $y$ conditioned on input $x$. This combination implicitly teaches the model to actively evaluate its current response at inference time and generate a self-revision if necessary. As a result, the trained model often exhibits explicit self-revision behaviors with a larger token cost at generation (see \Cref{sec:srt-results} and \Cref{sec:distillation-results} for a detailed discussion). In the next phase, we perform on-policy self-distillation on the \sft{} model, distilling this learned revision behavior back into the generator to enable sample-efficient self-improvement.

% This boosts the model's performance in its revision role. However, we also train the model on its own response $y$ and the subsequent $P_r$, conditioned on input $x$. 


% \ap{make the loss clearer}\jiarui{We didn't define what $y'$ is. Inconsistent naming: $y_{revised}$ vs. $y'$ vs. $z$. In my opinion it shouldn't be named as $y_{revised}$ because $\mathcal{D}_{\textsc{rephrase}}$ also uses it.} 

% \jiarui{I still don't understand this loss objective. Does this mean $-\sum_{t=1}^{|P_{y,r}| + |y_{\text{revised}}|}
% \log \pi_\theta\left(y_t \mid x, y_{<t}\right)$ for $\mathcal{D}_{\textsc{revision}}$, while $-\sum_{t=1}^{|y_{\text{revised}}|}
% \log \pi_\theta\left(y_t \mid x, P_{y,r}, y_{<t}\right)$ for $\mathcal{D}_{\textsc{rephrase}}$?}

% \sk{Alt proposal for notation: 
% \begin{itemize}
%     \item $a$: ground truth final answer (short string)
%     \item $y_{gold}$: correct first attempt response sampled from some model (could be eg. deepseek responses from the orig dataset); this could be useful for SFT or OPSD notation
%     \item $y$: initial reasoning response, sampled from student $\pi_{\theta (\cdot \mid x)}$
%     \item $c_r$: where $c_0$ = ``Wait, this is not correct, let me start over'' and $c_1$ = ``Let me rephrase the above solution.''
%     \item $z$: revised/rephrased response, sampled from  $\pi_{\theta (\cdot \mid x, y, c_r)}$
%     \item $\mathcal{D}_{revision} = \{ (x,y,z) : r(y, a) = 0, r(z, a) = 1\}$
%     \item $\mathcal{D}_{rephrase} = \{ (x,y,z) : r(y, a) = 1, r(z, a) = 1\}$
%     \item $\mathcal{L}_{\text{SRT}}({\theta})$ = 
% \end{itemize}
% }

% \yj{This loss is extremely confusing. Also, it's hard to know what is $y'$}\edward{I wonder if we even need to write out explicit NLL terms for this part.}


% We put more discussion in \ap{refer to the right section} on the importance of the two terms.\yh{also introduce rollout budget in appendix}
% \jiarui{In general, reviewers may be concerned that Phase 1 relies too heavily on the model’s self-revision capability?}

% \sk{it might be nice to add SCoRE baseline too (maybe after initial paper ddl if there's not enough time) ... SCoRE is basically like dapo/grpo but the model is given the opportunity to revise the initial response / try again}

% There are two terms in this loss function. The first term trains the model to play the role of the 'reviser', i.e. to revise or rephrase its response based on the initial response and the corresponding reward. The second term trains the model to retain its ability to generate good responses as a 'generator', by training on its correct initial responses in $ \mathcal{D}_{\textsc{rephrase}}$.


% This yields a self-generated critique dataset in which the supervision target is a newly generated correct reasoning trace rather than a gold chain-of-thought.

% \paragraph{Warmup objective.}
% We then optimize the model with a standard autoregressive objective on the revised responses:
% \[
% \mathcal{L}_{\mathrm{warm}}(\theta)
% =
% \mathbb{E}_{(p, z) \sim \mathcal{D}_{\mathrm{warm}}}
% \left[
% -\sum_{t=1}^{|z|}
% \log \pi_\theta \bigl(z_t \mid p, z_{<t}\bigr)
% \right].
% \]
% This warmup stage teaches the model an outcome-conditioned self-revision behavior: given its own sampled attempt and the ground-truth final answer, it learns to either rephrase a correct solution or restart from an incorrect one to produce a new correct reasoning response.

% \paragraph{On-policy response generation.}
% For each input $x$, we sample $K$ reasoning responses
% \[
% y^{(1)}, \dots, y^{(K)} \sim \pi_\theta(\cdot \mid x).
% \]
% We use a binary verifier $r(y, a) \in \{0,1\}$ to determine whether the final answer extracted from $y$ matches the ground-truth answer $a$.

% \paragraph{Outcome-conditioned feedback.}
% For each sampled response $y$, we construct a transformed supervision target $z$ as follows:
% \[
% z =
% \begin{cases}
% \rephrase(x, y), & r(y,a)=1, \\
% \refine(x, y, a), & r(y,a)=0.
% \end{cases}
% \]
% For correct responses, $\rephrase$ prompts the model $\pi_{\theta}$ to generate a rephrased version. For incorrect responses, $\refine$ prompts the model to generate a critique followed by a refined response conditioned on the ground-truth final answer. 

% In order to either perform $\rephrase$ or $\refine$, we use a critique prompt \critiqueprompt{}. To create our critique dataset, we remove all samples $(x, y, \critiqueprompt{}, z)$ for which reward $r(z,a)=0$,
% and create a critique dataset $\mathcal{D}_{\textsc{revision}} = \{(x, y, \critiqueprompt{}, z)\}$. 
% %The additional prompt that we use is given by \critiqueprompt{}.

%\yh{there's actually 2 phases with different loss.. How to unify them? Probably relate this to the twofold property}




% The second term computes loss on the generated response $z$, conditioned on input, output and the additional critique prompt: $(\inp, \out, \critiqueprompt{})$.



% This stage teaches the model an outcome-conditioned self-revision behavior: given its own sampled attempt and the ground-truth final answer, it learns to either rephrase a correct solution or restart from an incorrect one to produce a new correct reasoning response.


% This converts sparse outcome supervision into denser process-level supervision without requiring access to ground-truth reasoning traces.



% Idea:
%  - Student plays 2 roles: training and critique feedback
%  - Student generates responses, and the loss is computed using the logits of the teacher feedback

% If $\theta_{\textsc{revision}}$ denotes the student model after critique training,
% Self-revision training produces a self-revise reasoner model \sft{}, equipped with two-fold capabilities: 
% (1) \textbf{Self-Revision:} the model better critiques its own responses and comes up with an improved response.
% (2) \textbf{Reasoning:} the model generates better reasoning with self-revision behaviors.
%\yh{intuitively explain}
% More formally, define \textbf{student policy} $\pi_{\theta_{\student{}}} := \pi_{\theta^*}(\cdot \mid x)$. For any sampled response $(\inp, \out)$ from student policy, we obtain a binary outcome reward $r=r(y,a)\in \{0,1\}$ where $a$ denotes the ground truth answer. The \textbf{teacher policy} is then defined as the self-revise model's revised solution attempt, given the student's reasoning $y$ and a outcome reward $r$:
% \begin{equation}
% \pi_{\theta_{\teacher{}}} := \pi_{\theta^*} (y \mid P_r).
% \end{equation}

%\sk{(snippet below serves to transition phase 1 --> phase 2)}\\
% \sk{In summary, Phase 1 trains the model to revise incorrect responses. But at inference, we want correct first attempts, not just revision chains. This motivates Phase 2 of our pipeline.}


\subsection{Phase 2 (\opsd{}): On-Policy Self-Distillation via Revision Feedback}\label{sec:phase2}

%\yh{put motivation here}
\opsd{} phase distills the reviser’s behavior back into the generator, so that the model can internalize revision and produce more compact and well-directed responses. Concretely, we leverage the self-revision capability learned in Phase 1 to perform on-policy self-distillation. Let $\theta_{\sft{}}$ denote the model obtained at the end of Phase 1. In Phase 2, we initialize the student parameters as $\theta:=\theta_{\sft{}}$. 
The \textbf{generator} (student) uses the current parameters $\theta$ to produce on-policy responses $\pi_{\theta}(\cdot \mid x)$; its parameters are updated throughout Phase 2. The \textbf{reviser} (teacher) remains frozen at $\theta_{\sft{}}$ and provides token-level supervision by conditioning on the generator's response and its binary reward via $\pi_{\theta_{\sft{}}}(\cdot \mid x, y, P_r)$. The generator is trained to match the reviser's distribution with:
%$\theta:=\theta_{\sft{}}$ as the initial model checkpoint that plays both a generator and a reviser role: the generator role serves as the \textit{student} that outputs on-policy responses, while the reviser role serves as the \textit{teacher} that provides token-level supervision over its own responses. The resulting training objective is:
%\yh{Reverse the KL here: KL($\pi_s$ | $\pi_T$)}
% \sk{in the code, are you overriding $\alpha=0$ with $\alpha=1$? my understanding is that the sdft continual learning default code sets $\alpha=0$  (forward KL)}
% \begin{align*}
% \mathcal{L}_{\text{\opsd{}}}(\theta)
% &=
% \mathbb{E}_{(x,a)\sim\mathcal{D}}
% \;
% \mathbb{E}_{y\sim\pi_\theta(\cdot\mid x)}
% \sum_{t=1}^{|y|}
% \KL\!\left(
% \pi_\theta(\cdot \mid x, y_{<t})\;\|\;
% \pi_{\theta_{\sft{}}}(\cdot \mid x, y, P_{r}, y_{<t})\right).
% % \\
% % &\approx \mathbb{E}_{(x,a)\sim\mathcal{D}}
% % \;
% % \mathbb{E}_{y\sim\pi_\theta(\cdot\mid x)}
% % \sum_{t=1}^{|y|} \log \pi_\theta(y_t \mid x, y_{<t}) - \log \pi_{\theta_{\sft{}}}(\cdot \mid x, y, P_{r}, y_{<t})
% \end{align*}
{\small
\begin{equation}
\mathcal{L}_{\text{\opsd{}}}(\theta)
=
\mathbb{E}_{(x,a)\sim\mathcal{D}}
\;
\mathbb{E}_{y\sim\pi_\theta(\cdot\mid x)}
\sum_{t=1}^{|y|}
\KL\!\left(
\pi_\theta(\cdot \mid x, y_{<t})\;\|\;
\pi_{\theta_{\sft{}}}(\cdot \mid x, y, P_{r}, y_{<t})\right).
\label{eq:opsd}
\end{equation}
}
%\yh{put a simplified version of algorithm here}\\
% \sk{``in words'' alg summary below:} \\

\looseness-1 For each training example, a response is sampled from the model and checked against the ground-truth final answer for correctness. The Phase 1 reviser, conditioned on the model's response and whether it was correct, produces a next-token distribution at each token position. The model is trained to match this distribution via KL divergence loss.

% \looseness-1 In our main experiments, we fix the teacher model at the initial checkpoint and perform gradient updates only on the student to speed up training. We further show in \Cref{sec:self-evolution} that regularly synchronizing the teacher model with the updated student helps overcome training saturation and enables iterative self-evolution.

% \looseness-1 Through ablation studies, we find that the reviser’s distribution typically differs from the generator’s only at a sparse set of token positions that must be changed to reach a correct solution. Because the student’s response is generated on-policy, this supervision remains aligned with its current failure modes. As a result, Phase 2 compresses the model's verbose self-revision behaviors learned in Phase 1 into a compact, well-directed response.

% \sk{In other words, the teacher observes the student's complete response and corresponding binary reward, and --- having been trained in Phase 1 to revise --- converts this outcome-level signal into a per-token distributional target. The KL objective trains the student to match this target, with the learning signal concentrated at tokens where the student's reasoning diverges from a correct path. Because generation is on-policy for the student, this signal stays calibrated to the student's current weaknesses, and the net effect is that the two-step revision behavior from Phase 1 is compressed into single-pass generation via Phase 2.}

%This is equivalent to on-policy distillation and can be shown to be:
% \begin{align}
%    &\loss_{\ours{}}(\theta) = \kl ( \pi_{\theta_{\teacher{}}} || \pi_{\theta_{\student{}}} )
% \nonumber\\&=
% \mathbb{E}_{(x, a)} \mathbb{E}_{y\sim \pi_{\theta} (\cdot \mid x)}
% \sum_{t=1}^{|y|}   
% \left[
% \log \pi_{\theta^*} (y_t \mid x, y_{<t})
% - \log \pi_{\theta^*}(y_t \mid x, y_{<t})
% \right].
% \end{align}



% where $\stopgrad$ represents the Stop Gradient Operator.





% \begin{algorithm}[htbp]
% \caption{\ours{}}
% \label{alg:main}
% \textbf{Input:} Language model $\pi_\theta$; Dataset with questions $x$

% \begin{algorithmic}
% \State \Return 
% \end{algorithmic}
% \end{algorithm}



